\documentclass[12pt]{article}
\usepackage[utf8]{inputenc}
\usepackage[T1]{fontenc}
\usepackage[margin=1in]{geometry}
\usepackage{setspace}
\usepackage{xcolor}
\usepackage[colorlinks=true, urlcolor=blue!60!black, linkcolor=blue!60!black]{hyperref}
\setlength{\parindent}{0em}
\setlength{\parskip}{0.45em}
\renewcommand{\baselinestretch}{1.0}
\begin{document}
\begin{center}
{\Large \textit{Trump v. CASA, Inc.}, 606 U.S. 831 (2025)\par}
\vspace{0.3em}
{\normalsize Case Brief --- One-Page Edition\par}
\end{center}

\textbf{Memory Jogger.} Universal injunctions likely exceed federal courts' equitable power under the Judiciary Act (\href{https://supreme.justia.com/cases/federal/us/606/24a884/}{Justia}).

\textbf{Facts.} Executive Order No.~14160 denied citizenship documents to many U.S.-born children of noncitizens. Individuals, organizations such as CASA, and 22 States sued in three districts, alleging violations of the Citizenship Clause and the Nationality Act of 1940.

\textbf{Procedural History.} Three universal injunctions issued; the Government sought only partial stays.

\textbf{Votes.} Partial stays granted 6--3: Barrett for the Court; three concurrences; Sotomayor and Jackson dissented.

\textbf{Holding.} Relief is limited to what gives each plaintiff with standing complete relief; the order's legality was not decided.

\textbf{Reasoning.} Founding-era equity gave party-specific relief; universal injunctions are a twentieth-century innovation that evades Rule 23. Indivisible relief may incidentally help others; state plaintiffs' needs were remanded. Sotomayor and Jackson warned that rights now go unprotected for nonparties. \textit{Trump v. Barbara} (2026) later struck the order down through a class action.

\textbf{Examples.} \emph{Same-Side:} narrowing two plaintiffs' nationwide injunction. \emph{Opposite-Side:} a certified class. \emph{Fence-Sitter:} a State's cross-border claim.

\textbf{Critique.} Critics question freezing equity at 1789 practice, and classes, state suits, and APA vacatur have largely replicated nationwide relief.

\textbf{Key Quotation.} ``\,`Complete relief' is not synonymous with `universal relief.' It is a narrower concept: The equitable tradition has long embraced the rule that courts generally `may administer complete relief between the parties.'\,'' --- Justice Barrett.

\end{document}
